\documentclass[a4paper, serif, 10pt]{moderncv}

%% moderncv
% color and style
\moderncvcolor{black}
\moderncvstyle{banking}

%% page layout/margins
\usepackage[top=2.5cm, bottom=2.5cm, left=1.5cm, right=1.5cm]{geometry}

\nopagenumbers{}

%% Babel config for French language
% ref:
% - https://latex3.github.io/babel/guides/locale-french.html
\usepackage[french]{babel}

%% fontspec
\usepackage{fontspec}

\IfFontExistsTF{Linux Libertine}{
  \setmainfont[Ligatures={TeX, Common}, Numbers=OldStyle]{Linux Libertine}
}{}
\IfFontExistsTF{Linux Libertine O}{
  \setmainfont[Ligatures={TeX, Common}, Numbers=OldStyle]{Linux Libertine O}
}{}

%% CTeX
% CJK support, used to show CJK characters in the resume
%
% - heading=false: disable ctex heading
% - punct=kaiming: use kaiming punctuations styles for CJK
% - scheme=plain: use plain scheme, do not override \normalsize font size
% - space=auto: space settings for CJK characters
%
% CJK fonts are resolved per locale language, e.g. Noto Serif/Sans CJK JP
% for Japanese resumes, with fallback to the next region if not installed
%
% ref:
% - http://ctan.mirrorcatalogs.com/language/chinese/ctex/ctex.pdf
\usepackage[UTF8, heading=false, punct=kaiming, scheme=plain, space=auto]{ctex}

\IfFontExistsTF{Noto Serif CJK SC}{
  \setCJKmainfont{Noto Serif CJK SC}
}{}
\IfFontExistsTF{Noto Sans CJK SC}{
  \setCJKsansfont{Noto Sans CJK SC}
}{}

%% line spacing
\usepackage{setspace}
\setstretch{1.125}

%% URL styling - use normal text instead of monospace
\urlstyle{same}

%% Patch \httplink and \httpslink to handle full URLs
% The original moderncv macros always prepend http:// or https:// to URLs.
% This causes issues when the URL already has a protocol (e.g., https://example.com)
% resulting in malformed URLs like https://https://example.com.
% This patch checks if the URL already contains "://" and uses it directly if so.
% Note: We use \str_set:Nx (x-expansion) to expand macros like \@homepage first.
\ExplSyntaxOn
\str_new:N \g_yamlresume_url_str

\RenewDocumentCommand{\httpslink}{O{}m}{
  \str_set:Nx \g_yamlresume_url_str {#2}
  \str_if_in:NnTF \g_yamlresume_url_str {://}
    { \url{#2} }
    { \url{https://#2} }
}

\RenewDocumentCommand{\httplink}{O{}m}{
  \str_set:Nx \g_yamlresume_url_str {#2}
  \str_if_in:NnTF \g_yamlresume_url_str {://}
    { \url{#2} }
    { \url{http://#2} }
}
\ExplSyntaxOff

\name{Claire Moreau}{}
\title{Ingénieur en Intelligence Artificielle | Apprentissage profond et MLOps}
\phone[mobile]{+33 6 45 78 12 34}
\email{claire.moreau@example.com}
\homepage{https://clairemoreau.ai}

\address{24 rue de l'Innovation, Paris, Île-de-France, France, 75011}{}{}

\extrainfo{{\small \faLinkedin}\ \href{https://www.linkedin.com/in/clairemoreau-ai}{@clairemoreau-ai} {} {} {} • {} {} {} 
{\small \faGithub}\ \href{https://github.com/clairemoreau}{@clairemoreau} {} {} {} • {} {} {} 
{\small \faMedium}\ \href{https://medium.com/@clairemoreau}{@clairemoreau}}

\begin{document}

\maketitle

\section{Informations de base}

\cvline{}{Ingénieur en intelligence artificielle avec plus de 7 ans d'expérience dans la conception, l'entraînement et le déploiement de modèles d'apprentissage automatique à grande échelle. Spécialisée en vision par ordinateur et traitement du langage naturel, avec une expertise solide en MLOps et en industrialisation de solutions IA.
%
\begin{itemize}
\item Conception de pipelines de données et d'entraînement reproductibles avec Python, PyTorch et TensorFlow
\item Déploiement de modèles en production sur AWS et GCP avec Docker, Kubernetes et MLflow
\item Encadrement d'équipes pluridisciplinaires et vulgarisation des concepts IA auprès des parties prenantes
\item Veille technologique active et contribution à des projets open source
\end{itemize}}
\section{Formation}

\cventry{sept. 2017–juil. 2019}
        {Master, Intelligence Artificielle et Science des Données, Score : 16.5/20}
        {Université Paris-Saclay}
        {\url{https://www.universite-paris-saclay.fr}}
        {}
        {\begin{itemize}
\item Formation spécialisée en apprentissage profond, vision artificielle et traitement du langage
\item Projet de fin d'études sur la détection d'anomalies industrielles par autoencodeurs
\item Stage de recherche au sein de l'équipe IA de l'université
\end{itemize}
\textbf{Cours} : Apprentissage automatique, Vision par ordinateur, Traitement du langage naturel, Optimisation convexe, Statistiques bayésiennes, Systèmes distribués, MLOps et cloud computing}

\cventry{sept. 2014–juin 2017}
        {Licence, Mathématiques et Informatique, Score : 15.8/20}
        {Université de Bordeaux}
        {\url{https://www.u-bordeaux.fr}}
        {}
        {\begin{itemize}
\item Solide base en algorithmique, probabilités et statistiques
\item Initiation à la programmation scientifique et aux bases de données
\end{itemize}}
\section{Expérience professionnelle}

\cventry{janv. 2022–Present}
        {Ingénieur IA Senior}
        {NovaMind AI}
        {\url{https://www.novamind.ai}}
        {}
        {\begin{itemize}
\item Direction technique de la plateforme de vision par ordinateur pour le contrôle qualité industriel
\item Conception et entraînement de modèles de détection d'anomalies avec PyTorch et TensorFlow
\item Mise en place de pipelines MLOps reproductibles sur AWS SageMaker et Kubernetes
\item Encadrement de trois ingénieurs juniors et animation de revues de code
\item Collaboration avec les équipes produit et métier pour prioriser les cas d'usage à fort impact
\end{itemize}
\textbf{Mots-clés} : Apprentissage profond, Vision par ordinateur, MLOps, Python}

\cventry{sept. 2019–déc. 2021}
        {Ingénieur en Apprentissage Automatique}
        {DataForge Conseil}
        {\url{https://www.dataforge-conseil.fr}}
        {}
        {\begin{itemize}
\item Développement de modèles de traitement du langage naturel pour l'analyse de documents clients
\item Industrialisation de solutions de scoring et de recommandation avec Scikit-learn et Docker
\item Automatisation des workflows de données avec Airflow et Spark sur Google Cloud Platform
\item Rédaction de documentation technique et accompagnement des équipes métier dans l'adoption des modèles
\end{itemize}
\textbf{Mots-clés} : NLP, Scikit-learn, Cloud, Docker}
\section{Langues}

\cvline{Français}{Niveau natif ou bilingue}
\cvline{Anglais}{Niveau professionnel complet \hfill \textbf{Mots-clés} : TOEIC 950, Présentations techniques}
\cvline{Espagnol}{Niveau de travail limité \hfill \textbf{Mots-clés} : Espagnol des affaires}
\section{Compétences}

\cvline{Apprentissage automatique}{Expert \hfill \textbf{Mots-clés} : Python, PyTorch, TensorFlow, Scikit-learn, XGBoost}
\cvline{Traitement du langage naturel}{Avancé \hfill \textbf{Mots-clés} : Transformers, Hugging Face, spaCy, RAG, LangChain}
\cvline{MLOps et Cloud}{Avancé \hfill \textbf{Mots-clés} : Docker, Kubernetes, AWS SageMaker, MLflow, CI/CD}
\cvline{Science des données}{Avancé \hfill \textbf{Mots-clés} : Pandas, NumPy, SQL, Spark, Airflow}
\section{Récompenses}

\cventry{nov. 2023}
        {NovaMind AI}
        {Prix de l'innovation en IA}
        {}
        {}
        {Récompense décernée pour la conception d'un système de détection d'anomalies industrielles réduisant les faux positifs de 35 \%.}

\cventry{juil. 2019}
        {Université Paris-Saclay}
        {Meilleur mémoire de Master}
        {}
        {}
        {Distinction attribuée pour un mémoire sur les autoencodeurs variationnels appliqués à la maintenance prédictive.}
\section{Certificats}

\cventry{mars 2022}
        {AWS}
        {AWS Certified Machine Learning – Specialty}
        {\url{https://aws.amazon.com/certification/certified-machine-learning-specialty/}}
        {}
        {}

\cventry{juin 2021}
        {Google}
        {TensorFlow Developer Certificate}
        {\url{https://www.tensorflow.org/certificate}}
        {}
        {}
\section{Publications}

\cventry{oct. 2023}
        {Vers des modèles de vision robustes pour la détection d'anomalies industrielles}
        {Revue Française d'Intelligence Artificielle}
        {\url{https://www.rfia.fr/article/modeles-vision-robustes}}
        {}
        {\begin{itemize}
\item Étude de la robustesse des modèles de vision face aux variations d'éclairage et de capteurs
\item Évaluation comparative de plusieurs architectures convolutives sur des jeux de données industriels
\item Proposition d'une stratégie d'augmentation de données adaptée aux environnements de production
\end{itemize}}
\section{Références}

\cventry{\emaillink[sophie.bernard@novamind.ai]{sophie.bernard@novamind.ai}}
        {Directrice de l'ingénierie IA}
        {Sophie Bernard}
        {+33 1 23 45 67 89}
        {}
        {Claire allie une excellente maîtrise technique à une vraie capacité de leadership. Elle a su structurer notre plateforme MLOps et faire monter en compétence toute l'équipe.}
\section{Projets}

\cventry{févr. 2023–août 2023}
        {Plateforme de recherche augmentée par génération pour documents techniques}
        {Assistant documentaire RAG}
        {\url{https://github.com/clairemoreau/rag-assistant}}
        {}
        {\begin{itemize}
\item Indexation vectorielle de documents techniques avec FAISS et embeddings spécialisés
\item Chaîne de génération augmentée avec LangChain et un modèle LLM open source
\item Interface de question-réponse avec citations des sources et gestion des droits d'accès
\end{itemize}
\textbf{Mots-clés} : RAG, LangChain, LLM, FAISS}

\cventry{mars 2021–nov. 2021}
        {Système de vision industrielle pour la détection précoce de défauts}
        {Détection d'anomalies sur chaîne de production}
        {\url{https://github.com/clairemoreau/anomaly-detection}}
        {}
        {\begin{itemize}
\item Entraînement d'autoencodeurs convolutifs sur des images de pièces industrielles
\item Mise en place d'un tableau de bord de suivi des alertes et des performances du modèle
\item Réduction du taux de rebut de 12 \% lors du pilote sur une ligne de production
\end{itemize}
\textbf{Mots-clés} : Vision, Autoencodeurs, Industrie 4.0}
\section{Centres d'intérêt}

\cvline{Sport}{Escalade, Course à pied, Yoga}
\cvline{Culture}{Photographie, Lecture, Musique classique}
\cvline{Communauté}{Éthique de l'IA, Open Source}
\section{Bénévolat}

\cventry{sept. 2021–Present}
        {Mentore}
        {Women in AI France}
        {\url{https://www.womeninai.co/fr}}
        {}
        {\begin{itemize}
\item Accompagnement de femmes en reconversion vers les métiers de l'IA
\item Animation d'ateliers pratiques sur Python, la science des données et les entretiens techniques
\item Participation à des actions de sensibilisation dans des écoles et universités
\end{itemize}}
    
\end{document}